%!TEX program   = pdflatex
%!TEX encoding  = latin5 
%!TEX spellcheck = tr-TR 
%!BIB program    = bibtex

\documentclass[fleqn,oneside]{akdenizuni}
%%%%%%%%%%%%%%%%%%%%%%%%%%%%%%%%%%%%%%%%%%%%%%%%%%%%%%%%%%%%%%%%%%%%%%%%%%%%%%%%%%%%%%%%%%%%%%%%%%%%%%%%%%%%%%%%%%%%%%%%%%%%%%%%%%%%%%%%%%%%%%%%%%%%%%%%%%%%%%%%%%%%%%%%%%%%%%%%%%%%%%%%%%%%%%%%%%%%%%%%%%%%%%%%%%%%%%%%%%%%%%%%%%%%%%%%%%%%%%%%%%%%%%%%%%%%
\usepackage{makeidx}
\usepackage{amssymb}
\usepackage{eurosym}
\usepackage{graphicx}
\usepackage{amsmath}
\usepackage[pdfusetitle, pdfstartview=FitH,unicode]{hyperref}
\usepackage{amsthm}
\usepackage{times}
\usepackage[T1]{fontenc}
\usepackage{grffile}

\setcounter{MaxMatrixCols}{10}
%TCIDATA{OutputFilter=Latex.dll}
%TCIDATA{Version=5.50.0.2960}
%TCIDATA{<META NAME="SaveForMode" CONTENT="2">}
%TCIDATA{BibliographyScheme=BibTeX}
%TCIDATA{LastRevised=Thursday, February 07, 2019 15:22:59}
%TCIDATA{<META NAME="GraphicsSave" CONTENT="32">}
%TCIDATA{CSTFile=akdenizuni}

\theoremstyle{plain}
\newtheorem{theorem}{Teorem}[chapter]
\newtheorem{lemma}[theorem]{Lemma}
\newtheorem{problem}[theorem]{Problem}
\newtheorem{proposition}[theorem]{\"Onerme}
\newtheorem{remark}[theorem]{A\c c\i klama}
\newtheorem{corollary}[theorem]{Sonu\c c}
\newtheorem{claim}{\.{I}dda}
\theoremstyle{definition} 
\newtheorem{definition}[theorem]{Tan\i m}
\newtheorem{notation}[theorem]{G\"osterim}
\newtheorem*{solution}{\c C\"oz\"um}
\newtheorem*{example}{\"Ornek}
\newenvironment{keywords}{\labeledminipage[\keywordsname: ]}{\endlabeledminipage\par}
\newenvironment{onsoz}[1][\"ONS\"OZ]{\chapter*{#1}}{}
\renewenvironment{frame}[1][]{}{}
\providecommand\frametitle[2][]{}
\providecommand\framesubtitle[2][]{}
\allowdisplaybreaks
\input{tcilatex}
\def\part{\chapter}
\begin{document}

\title[KISA BA\c SLIK]{TEZ BA\c{S}LI\u{G}INI \"{U}\c{C} SATIRI GE\c{C}%
MEYECEK \c{S}EK\.{I}LDE BURAYA YAZINIZ}
\titleEN{ENGLISH TITLE NOT TO EXCEED THREE LINES SHOULD BE WRITTEN HERE}
\author[\.{I}. SOY\.{I}S\.{I}M]{\.{I}sim SOY\.{I}S\.{I}M}
\email{a@a.com}
\subject{Matemat\.\i k}
\subjectEN{Mathematics}
\danisman{Prof. Dr. \.{I}sim SOY\.{I}S\.{I}M}
\juri{Asst. Prof. Dr. \.{I}sim SOY\.{I}S\.{I}M}
\juri{Assoc. Prof. Dr. \.{I}sim SOY\.{I}S\.{I}M}
\ikincidanisman{Assoc. Prof. Dr. \.{I}sim SOY\.{I}S\.{I}M}
\university{T.C.\par
AKDEN\.{I}Z \"{U}N\.{I}VERS\.{I}TES\.{I}}
\institute{FEN B\.{I}L\.{I}MLER\.{I} ENST\.{I}T\"{U}S\"{U}}
\destek{Bu tez \dots\ taraf\i ndan \dots\ nolu proje ile desteklenmi\c{s}tir.}
%\date{26/10/2018}
\authorURL{http://matematik.fen.akdeniz.edu.tr/}
\maketitle

\begin{ozet}
\"{O}zet k\i sm\i\ tezin ilk etapta okunan k\i sm\i d\i r. Burada konunun
amac\i\ ve bu alanda \c{s}imdiye kadar yap\i lan \c{c}al\i \c{s}malar\i n
k\i saca neler oldu\u{g}undan yola \c{c}\i k\i larak kullan\i lan materyal
ve metot hakk\i nda bilgi verilmelidir. Yap\i lan deneyler ve \c{c}al\i \c{s}%
malardan \"{o}zet olarak bahis edilerek ula\c{s}\i lan sonu\c{c}lar
belirtilmelidir.

\.{I}yi yaz\i lm\i \c{s} bir \"{o}zet, ara\c{s}t\i rma yapan ki\c{s}inin
arad\i \u{g}\i\ bilgilerin bu \c{c}al\i \c{s}ma i\c{c}erisinde olup olmad\i 
\u{g}\i n\i n bilgisini tezin t\"{u}m\"{u}n\"{u} okumadan ara\c{s}t\i rmac\i
ya verebilmelidir. Ayr\i ca \c{c}al\i \c{s}man\i n konusu yan\i\ s\i ra
kapsam\i\ hakk\i nda da k\i sa s\"{u}rede fikir edinilmesine y\"{o}nelik
olmal\i d\i r. Buna ba\u{g}l\i\ olarak ara\c{s}t\i rmac\i , tezin genelini
inceleme karar\i n\i\ bu a\c{s}amada verebilmelidir. \"{O}zetin haz\i
rlanmas\i nda k\i sa ve \"{o}z c\"{u}mlelerle tez hakk\i nda ve elde edilen
sonu\c{c}lar y\"{o}n\"{u}nden ayd\i nlat\i c\i\ bilgiler bulunmal\i d\i r.
Bununla birlikte gereksiz detaya girmekten ka\c{c}\i n\i lmal\i d\i r.

\emph{\"{O}zet k\i sm\i\ i\c{c}in ortalama olarak yar\i m sayfa ile bir bu%
\c{c}uk sayfa aras\i nda bir a\c{c}\i klay\i c\i\ bilgi uygun olacakt\i r.
En az 150 kelime olmal\i d\i r.}

\begin{labeledminipage}[ANAHTAR KEL\.{I}MELER:]
Birinci anahtar kelime, \.{I}kinci anahtar kelime, \"{U}\c{c}\"{u}nc\"{u}
anahtar kelime, \dots\ \emph{ALFABET\.{I}K SIRALI OLMALIDIR!!!! ve EN AZ 4
ADET OLMALIDIR.}
\end{labeledminipage}
\end{ozet}

\begin{ozet}[english]
English abstract should be written here.

\begin{labeledminipage}[KEYWORDS:]
First keyword, Second keyword, Third keyword, \dots\ \ 
\end{labeledminipage}
\end{ozet}

\begin{onsoz}
\"{O}ns\"{o}z herhangi bir eser okunmadan \"{o}nce \c{c}al\i \c{s}may\i\
yapan ki\c{s}inin duygu ve d\"{u}\c{s}\"{u}ncelerini \c{c}al\i \c{s}madan ba%
\u{g}\i ms\i z olarak okuyuculara aktarabildi\u{g}i k\i s\i md\i r.

Burada eserin yazar\i , bu konunun neden se\c{c}ildi\u{g}ini, alan\i n \"{o}%
nemini, \c{c}al\i \c{s}may\i\ yaparken i\c{c}erisinde bulundu\u{g}u ko\c{s}%
ullar\i , bu esnada okuyucular\i\ ilgilendirebilece\u{g}ini d\"{u}\c{s}\"{u}%
nd\"{u}\u{g}\"{u} olaylar\i\ ve \c{c}al\i \c{s}may\i\ yaparken ya\c{s}ad\i
klar\i\ gibi konular\i\ belirtilebilir. \c{C}al\i \c{s}man\i n ba\c{s}%
kalar\i\ taraf\i ndan daha sonra geli\c{s}tirilmesi durumunda \"{o}nerilere
yer verilebilir.

Tezin haz\i rlanmas\i\ esnas\i nda yard\i m al\i nan ki\c{s}i ve kurum/kurulu%
\c{s}lara \"{o}ns\"{o}z i\c{c}erisinde bir paragraf ayr\i larak te\c{s}ekk%
\"{u}r edilebilir. Tezin \"{u}niversitemiz bilimsel ara\c{s}t\i rmalar proje
birimi taraf\i ndan desteklenmesi halinde bu durum i\c{c} kapakta
belirtilmelidir.
\end{onsoz}

\begin{glossary}[\protect\underline{\textbf{Simgeler:}}]
\item[\protect\underline{\textbf{Simgeler:}}] 

\item[A] Alfabenin ilk harfi

\item[B] Alfabenin ikinci harfi\newline

\item[\protect\underline{\textbf{K\i saltmalar:}} ] 

\item[MTA] Maden Tetkik Arama
\end{glossary}

\chapter{G\.{I}R\.{I}\c{S}}

Tezin ilk ve \"{o}nemli b\"{o}l\"{u}mlerinden birini olu\c{s}turan giri\c{s}
b\"{o}l\"{u}m\"{u} "G\.{I}R\.{I}\c{S}" ba\c{s}l\i \u{g}\i\ alt\i nda yaz\i
lmal\i\ ve en az iki sayfa olmal\i d\i r. Okuyucuya konuyu haz\i rlay\i c\i\
bilgiler verildikten sonra ara\c{s}t\i rman\i n amac\i\ ve kapsam\i , bu b%
\"{o}l\"{u}mde a\c{c}\i k\c{c}a belirtilmelidir. E\u{g}er tez \c{c}al\i \c{s}%
mas\i nda tart\i \c{s}mal\i\ bir adland\i rma, s\i n\i flama ve kavram
kullan\i lm\i \c{s}sa, bunlar\i n a\c{c}\i klamas\i\ yine "G\.{I}R\.{I}\c{S}%
" b\"{o}l\"{u}m\"{u}nde verilmelidir.

\chapter{KAYNAK TARAMALARI}

Tezin dayand\i \u{g}\i\ kuramsal bilgiler ve tez konusu ile ilgili olarak
belirtilmek istenen \"{o}nceki yay\i nlar, tezde tart\i \c{s}\i lan konular
ile ilgilendirilerek bu b\"{o}l\"{u}mde verilir ve en az iki sayfa olmal\i
d\i r.

\chapter{MATERYAL VE METOT}

Bu b\"{o}l\"{u}mde tezin materyal ve metoduyla ilgili bilgiler verilir.

\chapter{BULGULAR}

Bu b\"{o}l\"{u}mde tez \c{c}al\i \c{s}mas\i ndan elde edilen sonu\c{c}lar, 
\"{o}z ve a\c{c}\i k olarak yaz\i lmal\i d\i r.

\chapter{TARTI\c{S}MA}

Tez \c{c}al\i \c{s}mas\i ndan elde edilen bulgular\i n literat\"{u}rdeki
yeri \"{o}nceki \c{c}al\i \c{s}malar ve bulgular ile k\i yaslanarak ve
tart\i \c{s}\i larak bu b\"{o}l\"{u}m en az \"{u}\c{c} sayfadan olu\c{s}%
turulur. Tez konusunun niteli\u{g}ine ve ilgili Anabilim Dal\i\ kurallar\i
na g\"{o}re \textquotedblleft BULGULAR VE TARTI\c{S}MA\textquotedblright\
tek bir b\"{o}l\"{u}m alt\i nda da sunulabilir.

\chapter{SONU\c{C}LAR}

Bu b\"{o}l\"{u}mde, bu tez \c{c}al\i \c{s}mas\i ndan \c{c}\i kar\i lan sonu%
\c{c}lar, ilgili bilim dal\i n\i n temel ve uygulamal\i\ y\"{o}nlerine
yapt\i \u{g}\i\ ve yapabilece\u{g}i katk\i lar, bu \c{c}al\i \c{s}madaki
bulgulara dayanarak yap\i labilecek \"{o}neriler belirtilerek en az bir
sayfa olmal\i d\i r.

\chapter{KAYNAKLAR}

\nocite{*} 
\bibliographystyle{akdenizuni}
\bibliography{kaynaklar}

Tez i\c{c}eri\u{g}inde kullan\i lan kaynaklar, \emph{yazar soyad\i na g\"{o}%
re} s\i ralanmal\i d\i r.

Kaynaklar, a\c{s}a\u{g}\i daki genel kal\i ba uygun yaz\i lmal\i d\i r.
Yazar Soyad\i\ (b\"{u}y\"{u}k harf)-virg\"{u}l- \.{I}lk ad(lar)\i n\i n ba%
\c{s} harfi-nokta- yay\i n y\i l\i -nokta- Belgenin Ba\c{s}l\i \u{g}\i -
nokta- Yay\i nland\i \u{g}\i\ yer (yay\i n organ\i\ veya yay\i nevi)-virg%
\"{u}l -cilt no -virg\"{u}l- say\i\ no virg\"{u}l- sayfa no -virg\"{u}l-
Yay\i nland\i \u{g}\i\ \c{s}ehir veya \"{u}lke -nokta Ayn\i\ yazar(lar)\i n
farkl\i\ y\i llardaki yay\i nlar\i , tarih s\i ras\i na g\"{o}re
verilmelidir. Ayn\i\ yazar(lar)\i n ayn\i\ y\i ldaki yay\i nlar\i , yay\i n
y\i l\i n\i n hemen arkas\i nda a,b,c,\dots\ \c{s}eklinde belirtilerek
verilmelidir.

Kaynaklarda kullan\i lan periyodiklerin, uluslararas\i\ k\i saltmalar\i\
konmal\i d\i r. E\u{g}er bu k\i saltmalar bilinmiyorsa periyodik orijinal
ad\i yla yaz\i lmal\i d\i r. Periyodikler italik harflerle yaz\i lmal\i d\i
r.

\appendix

\chapter{EK1}

Metin i\c{c}inde yer almalar\i\ halinde konuyu da\u{g}\i t\i c\i\ ve okumada
s\"{u}reklili\u{g}i engelleyici nitelikte ve dip not olarak verilemeyecek
kadar uzun olan a\c{c}\i klamalar, bir form\"{u}l\"{u}n \c{c}\i kar\i l\i 
\c{s}\i , geni\c{s} kapsaml\i\ ve ayr\i nt\i l\i\ deney verileri, \"{o}rnek
hesaplamalar, vb. \emph{EKLER} b\"{o}l\"{u}m\"{u}nde verilmelidir. Bu b\"{o}l%
\"{u}mde yer alacak her bir a\c{c}\i klama i\c{c}in uygun bir ba\c{s}l\i k se%
\c{c}ilmeli ve bunlar sunu\c{s} s\i ras\i na g\"{o}re \emph{%
Ek-1,Ek-2,Ek-3,\ldots } \c{s}eklinde, her biri ayr\i\ bir sayfadan ba\c{s}%
layacak \c{s}ekilde sunulmal\i d\i r. EKLER \.{I}\c{c}indekiler dizininde
s\i ras\i yla verilmelidir.

\begin{ozgecmis}
\noindent  
\begin{minipage}{0.6\textwidth}
            \centering
            \makeatletter\@author\\\@email\makeatother 
            \end{minipage} \hfill  
\begin{minipage}{0.3\textwidth}% adapt widths of minipages to your needs
            \shorthandoff{=}
            \includegraphics[width=\linewidth]{foto}
            \shorthandon{=}
            \end{minipage}

\noindent 
\begin{tabular}{@{}ll}
\multicolumn{2}{l}{\textbf{\"{O}\u{G}REN\.{I}M B\.{I}LG\.{I}LER\.{I} }} \\ 
&  \\ 
Lisans: & Akdeniz \"{U}niversitesi \\ 
2012-2016 & Fen Fak\"{u}ltesi, Matematik B\"{o}l\"{u}m\"{u}, Antalya \\ 
&  \\ 
Y\"{u}ksek Lisans: & Akdeniz \"{U}niversitesi \\ 
2016-2019 & Fen Bilimleri Enstit\"{u}s\"{u}, Matematik Anabilim Dal\i ,
Antalya%
\end{tabular}
\end{ozgecmis}

\end{document}
